\section{偏差、方差与均方误差}
\label{sec:11-Mathematical-Statistics/02-Point-Estimation/01}

现在你手上至少有两个估计 \(\theta\) 的办法。对 Bernoulli 的 \(\theta\)，可以用样本均值 \(\bar X\)，也可以用别的——比如在样本量为零时不就变成先验猜测了吗？对正态均值，\(\bar X\) 是一个候选，样本中位数也是一个候选，或者收缩到零的 \(0.9\bar X\) 也可以考虑。

选哪个？需要一个评价体系——不是``这个看起来更好''，而是可以计算的、能比较的、在不同场景下仍然有意义的准则。

最直觉的想法：好的估计量应该``打中靶心''。重复抽样一百次，一百个估计值的平均值应该等于真值。这就是无偏性：
\[
\E_\theta[\hat\theta] = \theta.
\]
对所有 \(\theta\) 都成立。无偏估计量没有系统性偏移——它不会系统性地高估或低估任何一个参数值。这个要求很自然：如果每次估计都往同一个方向偏，那肯定有问题。

但只靠无偏性不够。看两个估计量：
\begin{itemize}
\tightlist
\item \(\hat\mu_1 = \bar X_n\)——样本均值。对正态数据，无偏，方差为 \(\sigma^2/n\)。
\item \(\hat\mu_2 = X_1\)——只用第一个观测。也是无偏（\(\E[X_1] = \mu\)），但方差是 \(\sigma^2\)，不随 \(n\) 减小。
\end{itemize}
两个都无偏。但第一个随着样本增加越来越精确，第二个不管你收集多少数据都和一个观测一样粗糙。无偏性区分不了它们。

更有甚者，无偏估计量可能给出不合法的参数值。Uniform\((0,\theta)\) 的矩估计 \(2\bar X\) 无偏（将在下一节推导），但可能小于样本最大值——一个认为部分已观测数据``不可能出现''的参数。有时候偏差不是敌人，产生非法估计值的无偏规则才是有问题的。

\subsection{均方误差把偏差和方差打包在一起}

评价估计量需要同时看两样：它瞄准哪里（偏差），以及它晃得多厉害（方差）。好的估计量偏差小、方差也小。但这两者经常冲突——可以通过``收缩''减小方差，代价是引入一点偏差。

在平方误差损失下，最常用的综合指标是均方误差（MSE）：
\[
\MSE_\theta(\hat\theta) = \E_\theta[(\hat\theta-\theta)^2].
\]
它衡量估计值与真值的平均平方距离。展开来看它由什么构成：
\[
\begin{aligned}
\hat\theta - \theta &= (\hat\theta - \E\hat\theta) + (\E\hat\theta - \theta), \\[4pt]
\MSE &= \E\bigl[(\hat\theta - \E\hat\theta)^2\bigr]
       + (\E\hat\theta - \theta)^2
       + 2(\E\hat\theta - \theta)\,\E[\hat\theta - \E\hat\theta].
\end{aligned}
\]
交叉项的期望为零：给定 \(\theta\)，\(\E\hat\theta - \theta\) 是常数，可以提出期望符号；剩下的 \(\E[\hat\theta - \E\hat\theta] = 0\) 直接来自均值的定义。因此：
\[
\boxed{\MSE = \Var_\theta(\hat\theta) + \Bias_\theta(\hat\theta)^2.}
\]

这个分解的意义很清楚：MSE 的两个来源，一个来自抽样波动（方差），一个来自系统偏移（偏差平方）。要降低 MSE，要么让估计更稳定（减方差），要么让估计更准确瞄准（减偏差），或者两者都做。但两者往往不能同时改善——这是统计估计中反复出现的交换。

\subsection{无偏不等于更准确}

设
\[
X_1,\ldots,X_n \sim N(\mu,\sigma^2),
\]
\(\sigma^2\) 已知。样本均值 \(\bar X\) 是无偏的，且在这个正态模型下是方差最小的无偏估计量。现在考虑一个有意引入偏差的收缩估计：
\[
\hat\mu_a = a\bar X,\qquad 0 \le a \le 1.
\]
当 \(a=1\) 时回到样本均值；\(a<1\) 时把估计往零拉——这是故意制造偏差。

逐项计算：
\[
\E[\hat\mu_a] = a\mu,\qquad \Bias = (a-1)\mu,
\]
\[
\Var(\hat\mu_a) = \frac{a^2\sigma^2}{n},
\]
因此：
\[
\MSE(\hat\mu_a) = (1-a)^2\mu^2 + \frac{a^2\sigma^2}{n}.
\]
对 \(a\) 求导并令其为零：
\[
-2(1-a)\mu^2 + \frac{2a\sigma^2}{n} = 0
\;\Longrightarrow\;
a^* = \frac{\mu^2}{\mu^2 + \sigma^2/n}.
\]

这个 \(a^*\) 说了什么？如果 \(|\mu|\) 很大——信号相对于噪声很强——\(a^*\approx 1\)，几乎不需要收缩，因为数据本身已经给出了清晰的信号。如果 \(|\mu|\) 很小——信号淹没在噪声中——\(a^*\approx 0\)，大幅收缩反而降低了 MSE，因为``零附近''平均比``被噪声污染的 \(\bar X\)''更靠近真 \(\mu\)。

但 \(a^*\) 依赖未知的 \(\mu\) 本身——你不能用真值来决定怎么估计真值。这是一个 oracle 结果，告诉你理论上存在比 \(\bar X\) 更优的估计量，而不是给你一个立即可用的公式。实际中，\(a\) 的选择通过先验、层次模型、交叉验证或 minimax 原则完成。

这个例子揭示的核心思想跨越了整个统计和机器学习：有意引入偏差来让估计更稳定，可能换取更小的总误差。正则化就是这个逻辑的系统化——岭回归的收缩、LASSO 的稀疏化、dropout 的随机扰动、早停的限制迭代——都是在某个调优参数控制下做偏差-方差交换。

\subsection{一致性讨论样本量增长后的行为}

MSE 评价的是固定样本量的表现。另一个角度是：随着 \(n\) 不断增长，估计量能不能``最终''逼近真值？

估计量序列 \(\hat\theta_n\) 一致，如果：
\[
\hat\theta_n \overset{P}{\longrightarrow} \theta,
\]
即对任意 \(\varepsilon > 0\)：
\[
P_\theta\bigl(|\hat\theta_n - \theta| > \varepsilon\bigr) \to 0 \quad (\text{当 } n\to\infty).
\]
Bernoulli 样本均值由（弱）大数定律一致。

一致性与无偏性互不等价：
\begin{itemize}
\tightlist
\item \(\bar X_n + 1/n\) 有偏（偏差 \(=1/n\)），但偏差随 \(n\) 趋零且方差也趋零，因此仍一致。小样本稍偏，大样本回到正轨。
\item \(\hat\mu_n = X_1\) 对所有 \(n\) 都无偏。但 \(\hat\mu_n - \mu \sim N(0,\sigma^2)\) 不随 \(n\) 收缩——偏离 \(\mu\) 超过任意 \(\varepsilon\) 的概率是固定的正数。意味着不管你收集多少额外数据后都扔掉不用，只用第一个观测——不一致。
\end{itemize}

一致性保证``足够多数据会解决这个问题''；但它不说``多少才算足够''，也不保证有限样本中表现好。反过来，一个在 \(n=30\) 时 MSE 优秀的规则，也可能在某种参数随样本量增长的方式下不一致。渐近结论和有限样本结论处于不同的逻辑层面，不能互相替代。

\subsection{风险函数把损失写进评价}

MSE 基于平方损失。更一般地，给定损失函数 \(L(\theta,a)\)，估计量 \(\delta(X)\) 的风险定义为：
\[
R(\theta,\delta) = \E_\theta\bigl[L(\theta,\delta(X))\bigr].
\]
平方损失产生 MSE。绝对误差损失 \(|\hat\theta-\theta|\) 关注中位表现，对极端偏差不如平方损失敏感。非对称损失——比如库存管理中低估需求比高估代价更大——对正负偏差分配不同权重，最优估计量会系统性地偏向代价更低的一侧。预测任务应评价新样本上的预测损失，而不是参数估计误差——两个任务的最优规则可以完全不同。

风险函数 \(R(\theta,\delta)\) 是参数 \(\theta\) 的函数。一个估计量可能在某些 \(\theta\) 值上风险小，在另一些上风险大。除非引入先验、参数范围约束或 minimax 等额外准则，通常不存在``对所有 \(\theta\) 都最好''的单一估计量。这就是为什么有这么多估计方法，各有各的擅长区域。

\subsection{标准误是方差的估计，不是总误差}

估计量的标准差：
\[
\SE_\theta(\hat\theta) = \sqrt{\Var_\theta(\hat\theta)}
\]
依赖未知参数，需要从数据估计。标准误量化的是抽样波动——如果重复抽样，估计值会在其期望周围晃多宽。它不包含偏差平方、模型错设和选择偏差。一个偏差巨大的估计量也可以有很小的标准误——每次都很稳定地偏到同一个错误方向。

报告``估计值 \(\pm\) 标准误''不是置信区间。置信区间还需要抽样分布形状（或渐近正态近似）以及覆盖概率的定义。对于偏斜分布、边界参数、或模型选择后的估计，正态对称的误差条可能严重误导——真实的抽样分布可能一边长一边短，或者整体偏离了名义中心。
